\def\ChapterCode{FORTBILDUNG}
	\chapter
[\CO{[FORTBILDUNG]} Ischämische Herzerkrankungen]
{Ischämische Herzerkrankungen}
\label{chp:ADVANCED-Atemweg}

\ChapterIntro
{%

}
{%
\begin{description}
  \item[Maintainer:] \ldots
  \item[Autoren:] Diverse
  \item[Reviewer:] Standard-Reviewprozess
  \item[Version:] {\VarDocVersionTypeName} ({\VarDocVersionTypeDescription})
  \item[SHA1:] {\DocGitCommit}
\end{description}

{\TxTeilkapitelAASS}}



\clearpage % TEMP temp clearpage


Das akute koronares Syndrom (ACS) ist ein Symptomenkomplex,
welcher typische Symptome ist chemischer Herzerkrankung
zusammenfasst: rassisch ist der Bremens-stechende
Thoraxschmerz, Atemnot, Anzeichen einer akuten
Herzinsuffizienz und EKG Veränderungen. Die Definition des
Krankheitsbildes ACS trägt der nur schwer möglichen
Unterscheidung von vorübergehenden und dauerhaften akuten
Störungen Rechnung. Eine Unterscheidung ob es sich um eine
vorübergehende Ischämie oder um eine tatsächliche in
Verzierung und dauerhafte Fluss störender Coronagefäße
handelt ist oft in der Akutsituation nur schwer oder gar
nicht zu treffen. Besonders bei einem unklaren EKG-Befund,
nicht eindeutigen Laborergebnissen oder nicht eindeutigen
Ergebnissen einer Echokardiographie kann nur durch eine
Verlaufsbeurteilung (serielle Laboruntersuchungen, serielle
EKG,…) Oder durch eine Angiographie der Nachweis oder der
Ausschluss einer relevanten dies chemischen Herzerkrankung
geführt werden.

\section{Vorspiel: Schlimmes bahnt sich an}

\section{Chronische Störungen: Koronare Herzkrankheit}

\section{Akute Störungen}

\subsection{Das ACS als Symptomenkomplex}

\subsection{Angina pectoris}

Angina Pectoris

Bei der Angina Pectoris kommt es zu einer vorübergehenden
Unterversorgung des Herzmuskels mit sauerstoffreichem Blut.
Ursächlich ist eine Durchflussstörung in einem oder mehreren
Herzkranzgefäße. Man unterscheidet je nach Klinik zwischen
einer stabilen und einer instabilen Angina Pectoris. Auch
bei einer instabilen Angina Pectoris kann es bereits zu
lebensbedrohlichen Komplikationen kommen.


\subsection{Myocardinfarkt}
Bei einem Myokardinfarkt kommt es zu einer Infarzierung,
d.\,h. ist chemischen Schädigung, infolge eines mangelnden
Flusses von sauerstoffreichem Blut zum Herzmuskel.
Ursächlich ist ein kompletter oder teilweiser, jedoch
wirksamer, Verschluss oder Verengung eines oder mehrerer
Herzkranzgefäße. Die Schwere des Myokardinfarkt ist abhängig
von der genauen Lokalisation der Durchblutungsstörung.
Kennzeichnend für die Blutversorgung des Herzmuskels ist die
kaum vorhandene Ausbildung von kulanter Aralkreisläufen, wie
dies anderswo im Körper üblich ist. Durch dieses Fehlen von
Kurt lateralen Kreisläufen bedeutet der Verschluss einer
Blutzuleitung nahezu automatisch eine Ischämie des
versorgten Gebietes. Je weiter am Anfang der koronarer
Gefäße der Verstoß liegt, desto größer die Auswirkungen auf
das versorgte Gebiet. Je weiter am Anfang der Coronagefäße
der Verschluss liegt, desto ein größeres Gebiet wird durch
diesen Verschluss in Mitleidenschaft gezogen. Je nach
Lokalisation des Verschlusses, d.\,h. Proxy mal oder distalen
Gefäß, betroffenes Gefäß) kann es zu unterschiedlichen
Symptomen kommen: besonders zeigt sich das deutlich im EKG,
wo die Auswirkungen auf die Betroffenen Ableitungen deutlich
sichtbar werden. Weiters können in der Elektrokardiogramm
Vieh Wandbewegungsstörung des Herzmuskels dargestellt und
die betroffene Region bzw. das vermutlich betroffene Gefäß
eingeschätzt werden. Erschwert wird dies jedoch dadurch,
dass die dass der Gefäßverlauf individuell unterschiedlich
ausgeprägt ist. So kann die Zuordnung eines betroffenen
Gebietes zu einem betroffenen Gefäß nur näherungsweise
erfolgen. Auch die Schmerzsymptomatik ist von der
betroffenen Region abhängig.
\subsection{Diagnostik}

\subsubsection{Leitsymptome, Anamnese und klinische Untersuchung}

\subsubsection{Labor}
Laboruntersuchungen sind ein wichtiger Teil der Diagnostik
der isländischen Herzerkrankung hierbei kommen verschiedene
Laborparameter zum Einsatz, welche unterschiedliche Aussage
Kraft haben. Gemeint ist diesem Laborparametern jedoch, dass
sie erst nach einer gewissen Zeit nach Beginn des
Ereignisses entsprechende Veränderungen erfahren, d.h. das
Laborparameter erst nach einer bestimmten Zeit positiv
ausfallen. Die wichtigsten Laborparameter für Diagnostik der
akuten bis chemischen Herzerkrankung sind die Herzenzyme.
Diese bestehen aus der Kreatinin King Nase (CK) sowie DIN
ISO Enzymen, der herzmuskelspezifischen Querdenker Nase
(CK-B). Weiters ist bei vorliegen einer systemischen
Situation die Lacktat die gedrückte Nase (LDH) erhöht. Der
wohl sind südlichste Parameter für das vorliegen einer ist
akuten bis chemischen Erkrankung ist das Druck Kundin, die
Zinses etwas ISO Enzym top ohnehin-T. Hierfür sind auch
Schnelltests zum Nachweis von top und NT, mit
eingeschränkter Sensitivität und Spezifität, verfügbar. Bei
unklaren ACS-Beschwerden kann die serielle Abnahme von den
erwähnten Laborparametern notwendig werden (zum Beispiel
alle drei, 6,12 Stunden).

\subsubsection{EKG}

Das EKG ist eine beliebte und weit verbreitete Möglichkeit,
Hinweise auf chemische Herzerkrankungen zu gewinnen.
Unmengen an Literatur beschäftigen sich mit der Zuordnung
von EKG Phänomenen zu ist chemischen Erkrankungen, sowie die
Interpretation derselbigen. Mit dem EKG ist es oft möglich,
den Ort der Störung oder die schwere der Störung sowie deren
Verlauf zu beurteilen und zu dokumentieren. Es bleibt jedoch
festzuhalten dass das EKG für sich alleine oft nur eine
begrenzte Aussagekraft hat: EKG Befunde sind jedenfalls
immer im Verlauf zu beurteilen, so kann auch ein
ursprünglich unauffälliges EKG innerhalb einer kurzen Zeit
bereits eindeutige Symptome einer Herzmuskelschädigung
zeigen. Ein EKG eignet sich jedenfalls nicht zum eindeutigen
Ausschluss einer chemischen Erkrankung.



\subsubsection{Echokardiographie}

Die ich hier Kartographie ist eine Ultraschalluntersuchung,
mit welcher die Wandbewegung des Herzmuskels, die Funktion
der Herzklappen unter Blutfluss innerhalb des Herzens
beurteilt werden kann. Die Methode ist nicht invasive und
kann unkompliziert durchgeführt werden. Ausschlaggebend für
eine korrekte gute Untersuchung ist jedoch eine gute Routine
des Untersucher. Klassische Befunde bei vorliegen einer ist
chemischen Herzerkrankung sind Wandbewegungsstörung infolge
einer fehlenden Funktion des Herzmuskels im betroffenen
Areal. Die Beurteilung, ob es sich um eine neu aufgetretene
oder um eine chronische Störung handelt, ist jedoch oft
nicht möglich. Einen großen Stellenwert hat die
Echokardiographie auch im Ausschluss von
Differentialdiagnosen wie zum Beispiel dem Aneurysma (Thora
kahl, Herz Aneurysma,…) Oder anderen Störungen, welche
ähnliche Symptome verursachen (Hinweise auf per Gremium
Kapitels, Rekap Ergüsse,…).tion{Diagnostische Angiographie}
Diagnostische Angiographie

Die Angiographie bietet eine ausgezeichnete Möglichkeit zur
Beurteilung von Coronagefäßen. Sie ist einerseits
diagnostisch durchführbar, andererseits können mittels
Angiographie auch Interventionen durchgeführt werden, wie
zum Beispiel Anlage eines Stent oder die Dehnung des
Herzkranzgefäße mittels eines Hallo muss. Nachteilig ist
jedoch war der relativ hohe Aufwand der notwendig ist, die
Untersuchung ist invasive und es sind schwer wiegende
Komplikationen möglich. Weiters verfügen viele Krankenhäuser
nicht über eine entsprechende Einrichtung bzw. entsprechend
geschultes Personal. Auch hier ist die Routine und das
können des Untersucher wesentlich für den Erfolg der
Untersuchung (und gegebenenfalls für den Erfolg der
Intervention).
\subsection{Management}

\subsubsection{Übersicht und Erstmaßnahmen}

\subsubsection{Observanz}

\subsubsection{Interventionelles Vorgehen}

\subsubsection{Lyse}

\subsubsection{Medikamentöse Begleittherapie}



\section{Folgestörungen}




\ChapterEnd